\chapter{أدوات العمل اليومي والتواصل}
\label{ch:daily}

تتعرف على التطبيقات الجاهزة التي تستخدمها كل يوم: من أين تبدأ يومك، وكيف تراسل زملاءك وتجتمع بهم وتحفظ ملفاتك، وكيف تتابع مهامك وموافقاتك.

\section{خريطة الأدوات الجاهزة}
\label{sec:daily-map}
\begin{rcslide}{خريطة الأدوات الجاهزة}
\sphase{نظرة عامة}
\slead{تجد في كيان من أول يوم تطبيقات جاهزة لكل عمل يتكرر في المؤسسات، وكلها تعمل على بيانات واحدة.}
\scard{chat}{التواصل}{البريد الإلكتروني، والمحادثات، والتقويم.}
\scard{tasks}{العمل اليومي}{المهام، وطلبات الموافقة، ومكتبة الملفات.}
\scard{building}{تطبيقات المؤسسة}{الموارد البشرية، والاستراتيجية والتنفيذ، والمراسلات، ودليل الإجراءات.}
\scard{lock}{الإعدادات}{القواعد التي تعمل بها المؤسسة كلها.}
\stake{تتعرف في هذه الوحدة على أدوات يومك، وفي الوحدة التالية على تطبيقات المؤسسة.}
\slink{main}{الدليل الرئيسي}{ar/index.html}
\end{rcslide}

\section{تبدأ يومك من الرئيسية}
\label{sec:daily-home}
\begin{rcslide}{تبدأ يومك من الرئيسية}
\sphase{يومك}
\slead{تجمع الصفحة الرئيسية كل ما ينتظرك من التطبيقات كلها، فلا تفتحها واحداً واحداً لتعرف ما عليك.}
\smetric{الموافقات}{تبحث عنها}{«موافقات بانتظارك»}
\smetric{المهام}{تتذكرها}{«مهام مستحقة»}
\smetric{الاجتماعات}{في تقويم منفصل}{«اجتماعات اليوم»}
\smetric{المتأخر}{تكتشفه متأخراً}{تنبيه في أعلى الصفحة}
\stake{افتح الرئيسية أول اليوم: إن كان شيء متأخراً أو ينتظر قرارك، فستراه هنا أولاً.}
\end{rcslide}

\section{البريد الإلكتروني}
\label{sec:daily-email}
\begin{rcslide}{البريد الإلكتروني}
\sphase{التواصل}
\slead{صندوق بريد كامل داخل كيان: ترسل وتستقبل وترد وتعيد التوجيه، بمرفقات ومجلدات وعلامات.}
\scard{mail}{بريد كيان أو حسابك}{صندوق مستضاف في كيان، أو حسابك في Gmail أو Outlook وغيرهما.}
\scard{paper}{المجلدات والعلامات}{الوارد والمرسل والمسودات والأرشيف، وعلاماتك الخاصة.}
\scard{ai}{لوحة الذكاء الاصطناعي}{ملخص الرسالة، وكشف الاحتيال والرسائل المزعجة، والمهام والإجراءات.}
\scard{tasks}{المهام المرتبطة}{تربط الرسالة بالمهمة التي تخصها.}
\stake{بريدك في المكان نفسه الذي فيه مهامك وتقويمك، فلا تنتقل بين برامج لتتابع عملاً واحداً.}
\end{rcslide}

\section{المحادثات}
\label{sec:daily-chat}
\begin{rcslide}{المحادثات}
\sphase{التواصل}
\slead{محادثات فورية مع زملائك، فردية أو في مجموعات، بدل مجموعات واتساب الخاصة بالعمل.}
\scompare{مجموعات واتساب للعمل}{تختلط بالرسائل الشخصية \\ الملفات تضيع في المحادثة \\ الطلب يبقى رسالة}{المحادثات في كيان}{للعمل وحده، ولأعضاء المؤسسة \\ ملفات ورسائل صوتية تُفرَّغ نصاً \\ «إنشاء مهمة» من الرسالة نفسها}
\scard{people}{مجموعات وقنوات}{تدعو إليها الأعضاء.}
\scard{chat}{الرد في سلسلة}{وتثبيت الرسائل المهمة.}
\scard{ai}{البوتات}{مساعدون يعملون معك في المحادثة.}
\stake{حين يتحول طلب في المحادثة إلى مهمة، يصبح له مسؤول وموعد، ولا يضيع بين الرسائل.}
\end{rcslide}

\section{التقويم والاجتماعات}
\label{sec:daily-calendar}
\begin{rcslide}{التقويم والاجتماعات}
\sphase{التواصل}
\slead{مواعيدك واجتماعاتك في مكان واحد، مع الحضور والقاعات وما دار في كل اجتماع.}
\scard{flow}{الأحداث والتكرار}{باليوم أو الأسبوع أو الشهر، وأوقات للتركيز.}
\scard{people}{الدعوات}{«هل ستحضر؟» نعم أو ربما أو لا، ومقارنة جداول الزملاء.}
\scard{building}{القاعات}{ينبّهك التقويم إذا كانت القاعة محجوزة.}
\scard{paper}{محضر الاجتماع}{جدول الأعمال والمحضر والخطوات التالية.}
\stake{الاجتماعات عبر الإنترنت تتم برابط خارجي تضعه في الحدث، فيظهر للحضور زر «انضمام».}
\end{rcslide}

\section{مكتبة الملفات}
\label{sec:daily-library}
\begin{rcslide}{مكتبة الملفات}
\sphase{العمل اليومي}
\slead{ملفات المؤسسة في مساحات ومجلدات مشتركة، بدل أن تبقى على أجهزة الموظفين.}
\scard{building}{المساحات المشتركة}{مجلدات للفريق أو للمؤسسة.}
\scard{link}{المشاركة}{«مشترك معي» و«شاركته أنا»، ونسخ الرابط.}
\scard{form}{المعاينة والتحرير}{تفتح المستند وتعدّله في المتصفح.}
\scard{flow}{سجل الإصدارات}{تعود إلى أي نسخة سابقة.}
\stake{لا نسخ متعددة من الملف نفسه بالبريد: ملف واحد، وكل تعديل فيه محفوظ في سجل إصداراته.}
\end{rcslide}

\section{المهام: من طلب عابر إلى مهمة واضحة}
\label{sec:daily-tasks}
\begin{rcslide}{المهام: من طلب عابر إلى مهمة واضحة}
\sphase{العمل اليومي}
\slead{كل عمل له مسؤول وموعد وحالة محدّثة، فيعرف الفريق ما أُنجز وما يعيق العمل.}
\scompare{قبل}{رسالة: «تابع هذا الموضوع» \\ لا موعد ولا مسؤول واضح \\ تسأل لتعرف أين وصل}{في كيان}{مهمة بعنوان محدد \\ لها مسؤول وموعد \\ حالتها وعوائقها مكتوبة فيها}
\scard{tasks}{المجموعات}{مشتركة لعمل الفريق، أو شخصية لعملك الخاص.}
\scard{chat}{التقدم والعوائق}{تسجّل ما أنجزته وما يعيقك، فيتابع زميلك دون تخمين.}
\stake{المهمة تحفظ تاريخ العمل كله في مكان واحد، فلا يضيع شيء بين المحادثات.}
\slink{main}{المجموعات}{ar/tasks-collections.html}
\slink{main}{إنشاء المهمة وتفاصيلها}{ar/tasks-details.html}
\slink{main}{التقدم والنتائج والتواصل}{ar/tasks-follow-up.html}
\end{rcslide}

\section{طلبات الموافقة}
\label{sec:daily-approvals}
\begin{rcslide}{طلبات الموافقة}
\sphase{العمل اليومي}
\slead{كل ما ينتظر قرارك في مكان واحد، أياً كان التطبيق الذي جاء منه.}
\scard{tasks}{قيد الانتظار}{ما ينتظر قرارك الآن.}
\scard{people}{أنشأتها أنت}{الطلبات التي قدّمتها، وأين وصلت.}
\scard{paper}{قرارك}{«موافقة» أو «رفض» أو «طلب تغييرات»، وقد يُطلب توقيعك.}
\scard{flow}{المراحل}{تعرف في أي مرحلة الطلب، ومن يأتي بعدك.}
\stake{إذا تأخر طلبك عند أحد، ترسل له تذكيراً من الطلب نفسه، بدل أن تلاحقه بالهاتف.}
\end{rcslide}

\section{أدوات تتحدث فيما بينها}
\label{sec:daily-linked}
\begin{rcslide}{أدوات تتحدث فيما بينها}
\sphase{الخلاصة}
\slead{لأن هذه الأدوات في منصة واحدة، ينتقل العمل من أداة إلى أخرى دون نسخ أو إعادة كتابة.}
\schain{رسالة في المحادثات ← مهمة لها مسؤول وموعد}
\schain{دعوة اجتماع في البريد ← حدث في التقويم}
\schain{رسالة بريد ← المهمة المرتبطة بها}
\stake{كل ما يصلك، أياً كان مصدره، يمكن أن يصبح عملاً واضحاً تتابعه حتى نهايته.}
\end{rcslide}
